\documentclass[10pt,a4paper]{amsart}
\usepackage[margin=19mm]{geometry}
\usepackage{amsmath,amssymb,mathtools}
\usepackage[scr=rsfs,cal=euler]{mathalfa}
\usepackage{xfrac}
\usepackage[unicode,hidelinks,bookmarks=false]{hyperref}
\newcommand{\ie}{i.e.\ }
\newcommand{\cf}{cf.\ }
\newcommand{\resp}{respectively\ }
\newcommand{\Subsection}{\subsection}
\providecommand{\nobreakdash}{\nobreak\hbox{-}}
\setlength{\emergencystretch}{2em}
\multlinegap0pt
\usepackage{bm}
\usepackage{bbm}
\usepackage{dsfont}
\usepackage{xspace}
\usepackage{cancel}
\usepackage{mathtools}
\usepackage{amsthm}
\makeatletter
\@ifundefined{thm}{\newtheorem{thm}{Thm}}{}
\@ifundefined{lemma}{\newtheorem{lemma}[thm]{Lemma}}{}
\@ifundefined{prop}{\newtheorem{prop}[thm]{Proposition}}{}
\makeatother
\def\cL{\mathcal{L}}
\def\dd{\mathrm{d}}
\newcommand{\escal}[1]{\langle#1\rangle}
\renewcommand{\ker}{\mathrm{Ker}}
\title[Kenmotsu manifolds and spin-c Killing spinors]{Kenmotsu manifolds and spin-c Killing spinors}
\author{Jos\'e Luis Carmona Jim\'enez, Alejandro Gil-Garc\'ia, C. S. Shahbazi}
\date{Source: arXiv:2608.19108v1 (CC BY 4.0)}
\begin{document}
\maketitle

\noindent\emph{Source and licence.} Kenmotsu manifolds and spin-c Killing spinors. Version arXiv:2608.19108v1 (CC BY 4.0), distributed under the Creative Commons Attribution 4.0 International licence, as recorded in the arXiv licence metadata for this exact version and shown on the abstract page https://arxiv.org/abs/2608.19108v1. Full rights notice: licenses/arxiv-2608.19108-CC-BY-4.0.txt.\par\smallskip

\noindent\textit{Editorial scope.} Counted unit: the Lemma whose source label is \texttt{lemma:nabla\_zeta} in the article (source0.tex lines 707--709, source label \texttt{lemma:nabla\_zeta}), delivered together with the article's own complete original proof of it (source0.tex lines 711--782, one \texttt{proof} environment of 72 lines). No mathematics is written, repaired or re-derived: statement and proof are the article's own text line for line. Required context retained uncounted: prop:KillingSpinorBilinear (lines 589--599). Every definition, display and label of those lines is retained unchanged. Omitted from this package and not redistributed: 1--588, 600--706, 710--710, 783--1282. Nothing inside a retained excerpt is omitted or altered: every retained body is delivered exactly as the article's lines are written, so no omission receipt of any kind accompanies this package. Citations: the retained excerpts carry no citation command at all, so no bibliography entry of the article is transcribed and no reference is redistributed. Reference adaptation: none was needed and none was made; every reference inside the delivered unit resolves inside it, so no unresolved reference or citation is produced. Printed slips kept literal: no printed word, symbol, hypothesis or number of the counted statement or of its proof was altered. Layout and class: the article's own class is \texttt{amsart} with packages including geometry, todonotes, bm, latexsym, amsmath, amstext, amssymb, amsfonts, amscd, array, multirow, amsbsy, mathrsfs, stmaryrd; the wrapper is the lane's pinned \texttt{amsart} editorial wrapper at 10pt on A4, which restates the theorem-like environments the retained text needs and adds the article's own macro definitions that the retained text needs (\texttt{\textbackslash cL}, \texttt{\textbackslash dd}, \texttt{\textbackslash escal}, \texttt{\textbackslash ker}), with extra wrapper packages (bm, bbm, dsfont, xspace, cancel, mathtools). The delivered numbering is the unit's own. Assets: no retained span contains a figure environment, a graphics-inclusion command, a PDF-page inclusion, a table or a TikZ picture; no image file is copied into this package, the package declares no asset dependency and no image file is redistributed with it. Licence: version arXiv:2608.19108v1 of this article is distributed under the Creative Commons Attribution 4.0 International licence, as recorded in the arXiv licence metadata for this exact version and shown on the abstract page https://arxiv.org/abs/2608.19108v1. A full rights notice is delivered at licenses/arxiv-2608.19108-CC-BY-4.0.txt. Characters: every retained excerpt is pure ASCII, so the delivered engine, XeLaTeX with its default Latin Modern text font, reports no missing character.
\begin{prop}
\label{prop:KillingSpinorBilinear}
A $(2n+1)$-dimensional oriented and spin-c Riemannian manifold $(M,g)$ admits a pure spin-c Killing spinor with imaginary Killing function $i\mu$, where $\mu \in C^\infty(M, \mathbb{R})$, if and only if it admits a Hermitian complex line bundle $\cL$ and a Cartan $n$-form $\rho \in \Omega^n_{\mathbb{C}}(M, \mathcal{L})$ satisfying the differential equation:
\begin{equation}
\label{eq:imaginary-killing-rho}
\nabla^{g,A}_X \rho = i^n \mu \ast (X^\flat \wedge \rho)
\end{equation}

\noindent
for every vector field $X \in \Gamma(TM)$.
\end{prop}

\begin{lemma}\label{lemma:nabla_zeta}
The one-form $\zeta\in\Omega^1(M)$ satisfies $\nabla^g_X\zeta = \mu(X^\flat - \zeta(X)\zeta)$ for all $X\in \Gamma(TM)$. In particular, $\zeta$ is closed and $\mathcal{D} := \ker(\zeta)$ is an integrable distribution.
\end{lemma}

\begin{proof}
Let $c_n:=i^n\mu$. Then, by Proposition~\ref{prop:KillingSpinorBilinear}, we have $\nabla^{g,A}_X\rho=c_n*(X^\flat\wedge\rho)$ for all $X\in\Gamma(TM)$. We first compute the differential of the squared norm $u:=\escal{\rho,\overline{\rho}}$:
\begin{equation*}
X(u)=c_n\escal{*(X^\flat\wedge\rho),\overline{\rho}}+\overline{c}_n\escal{\rho,*(X^\flat\wedge\overline{\rho})}.
\end{equation*}

\noindent
Note that:
\begin{equation*}
\escal{*(X^\flat\wedge\rho),\overline{\rho}}=\escal{X^\flat\wedge\rho,*\overline{\rho}}=\escal{\rho,\iota_X(*\overline{\rho})}=\escal{\rho,*(\overline{\rho}\wedge X^\flat)}=(-1)^n\escal{\rho,*(X^\flat\wedge\overline{\rho})}.
\end{equation*}

\noindent
Hence $X(u)=((-1)^nc_n+\overline{c}_n)\escal{\rho,*(X^\flat\wedge\overline{\rho})}$. Since $\mu$ is a real function, we have $\overline{c}_n=(-i)^n\mu=(-1)^n c_n$. The two terms thus add giving:
\begin{equation*}
X(u)=2\overline{c}_n\escal{\rho,*(X^\flat\wedge\overline{\rho})}.
\end{equation*}

\noindent
Using $i^n\rho\wedge\overline{\rho} = u *\zeta$, which implies $\rho\wedge\overline{\rho} = (-i)^n u *\zeta$, we have:
\begin{equation*}
\escal{\rho,*(X^\flat\wedge\overline{\rho})}\nu = \rho\wedge *(*(X^\flat\wedge\overline{\rho})) = (-1)^n X^\flat\wedge\rho\wedge\overline{\rho} = (-1)^n(-i)^n u\zeta(X)\nu = i^n u \zeta(X) \nu,
\end{equation*}

\noindent
where $\nu$ is the Riemannian volume form. Thus:
\begin{equation*}
X(u)=2(-i)^n\mu (i^n u\zeta(X))=2\mu u\zeta(X).
\end{equation*}

\noindent
Using this, we compute the covariant derivative of $\zeta$:
\begin{align*}
\nabla^{g}_X\zeta &= \nabla^{g}_X \left(\frac{i^n}{u} *(\rho\wedge\overline{\rho})\right)\\ 
&= -2\mu\zeta(X)\zeta+\frac{i^n}{u}*\big(c_n*(X^\flat\wedge\rho)\wedge\overline{\rho} + \overline{c}_n\rho\wedge *(X^\flat\wedge\overline{\rho})\big).
\end{align*}

\noindent
Now we evaluate the tensor contribution. Define the $(0,2)$-tensor field $T$ by $T(X,Y):=\iota_Y(*(*(X^\flat\wedge\rho)\wedge\overline{\rho}))$ for all $X$, $Y\in\Gamma(TM)$. Then:
\begin{align*}
    T(X,Y)&=\iota_Y(*(*(X^\flat\wedge\rho)\wedge\overline{\rho}))=*(*(X^\flat\wedge\rho)\wedge\overline{\rho}\wedge Y^\flat)\\
    &=(-1)^n*(Y^\flat\wedge\overline{\rho}\wedge*(X^\flat\wedge\rho))=(-1)^n*\escal{Y^\flat\wedge\overline{\rho},X^\flat\wedge\rho}\nu\\
    &=(-1)^n\escal{X^\flat\wedge\rho,Y^\flat\wedge\overline{\rho}}=(-1)^n\big(g(X,Y)u-\escal{\iota_X\overline{\rho},\iota_Y\rho}\big).
\end{align*}

\noindent
Observe that $\rho \wedge \ast(X^\flat \wedge \overline{\rho}) = (-1)^n \ast(X^\flat \wedge \overline{\rho}) \wedge \rho$. The bracketed term in the covariant derivative evaluated on $Y$ is therefore $c_n T(X,Y) + \overline{c}_n (-1)^n \overline{T(X,Y)}$. Since $\overline{c}_n=(-1)^nc_n$, this equals $c_n \big(T(X,Y) + \overline{T(X,Y)}\big)$. Multiplying by $\frac{i^n}{u}$, and noticing $i^n c_n = i^{2n}\mu = (-1)^n\mu$, we obtain:
\begin{equation*}
(-1)^n\frac{\mu}{u} \big(T(X,Y) + \overline{T(X,Y)}\big) = \frac{\mu}{u} \big(2g(X,Y)u - (\escal{\iota_X\overline{\rho},\iota_Y\rho} + \escal{\iota_X\rho,\iota_Y\overline{\rho}})\big).
\end{equation*}

\noindent
Since $\rho = (\theta^1\wedge \cdots \wedge \theta^n) \otimes \mathfrak{l}$ is a Cartan $n$-form built from isotropic and mutually orthogonal components, evaluating the symmetric sum over its factors yields exactly the metric on the distribution $\mathcal{D}=\ker(\zeta)$:
\begin{equation*}
\escal{\iota_X\overline{\rho},\iota_Y\rho} + \escal{\iota_X\rho,\iota_Y\overline{\rho}} = \frac{u}{2}\sum_{j=1}^n(\theta^j\odot\bar{\theta}^j)(X,Y) = u\big(g(X,Y)-\zeta(X)\zeta(Y)\big).
\end{equation*}

\noindent
Substituting this back, we obtain $\frac{\mu}{u} \big( u g(X,Y) + u \zeta(X)\zeta(Y) \big) = \mu \big( g(X,Y) + \zeta(X)\zeta(Y) \big)$. Subtracting the variation of the norm contribution $2\mu\zeta(X)\zeta(Y)$ yields:
\begin{equation*}
(\nabla^g_X\zeta)(Y) = \mu\big(g(X,Y)-\zeta(X)\zeta(Y)\big).
\end{equation*}

\noindent
Therefore, we obtain the formula in the statement. Since $\nabla^g\zeta$ is a symmetric tensor:
\begin{equation*}
\dd\zeta(X,Y)=(\nabla_X^g\zeta)(Y)-(\nabla_Y^g\zeta)(X)=0,
\end{equation*}

\noindent
thus $\zeta$ is closed and $\mathcal{D}=\ker(\zeta)$ is an integrable distribution by the Frobenius theorem.
\end{proof}

\end{document}
